\section{Stage 3: Boundary Extraction and Greeks}

\subsection{Purpose and Method}
Once the CN/PSOR surface is available, Stage 3 derives secondary quantities
without pretending that they are as numerically stable as price. For every time
level, compute
\[
  C_i^n=U_i^n-\Phi(S_i).
\]
Nodes with \(C_i^n\) at or below a small tolerance are classified as exercise
nodes; larger premiums indicate continuation.

\subsection{Boundary orientation and status}
For a put, the exercise set begins at low spot. The algorithm finds the first
continuation node and interpolates between it and the preceding exercise node.
For a dividend-paying call, exercise is on the high-spot side. The algorithm
finds the last continuation node and interpolates toward the next exercise
node. These directions are deliberately different.

The output records one of three statuses:
\begin{itemize}
  \item \texttt{boundary\_found}: both exercise and continuation nodes exist;
  \item \texttt{no\_exercise}: no exercise node appears on the truncated grid;
  \item \texttt{all\_exercise}: every grid node is classified as exercise.
\end{itemize}
If no boundary exists, the code records a missing value rather than fabricating
a threshold.

\begin{commonconfusion}
A no-dividend American call control should generally have no early-exercise
boundary. Replacing its missing boundary with zero or \(S_{\max}\) would turn a
valid theorem-based control into a false supervised label.
\end{commonconfusion}

\subsection{Delta, Gamma, and strict masks}
For interior nodes, the teaching approximations are
\[
  \Delta_i^n
  \approx
  \frac{U_{i+1}^n-U_{i-1}^n}{2\Delta S},
\]
\[
  \Gamma_i^n
  \approx
  \frac{U_{i+1}^n-2U_i^n+U_{i-1}^n}{(\Delta S)^2}.
\]
Delta is a first derivative; Gamma is a second derivative. Differentiation
amplifies grid and interpolation error, and the second derivative amplifies it
again. The payoff kink at \(S=K\) and the exercise transition at \(S=B(\tau)\)
are therefore delicate regions.

The strict Greek mask excludes nodes near zero spot, the payoff kink, and a
finite band around the free boundary. The mask does not make Gamma exact. It
defines a less singular region in which comparisons are more meaningful.

\begin{readerintuition}
Price can look smooth and accurate while Gamma oscillates, because many nearby
price errors cancel in the level but become magnified by repeated
differencing.
\end{readerintuition}

\section{Stage 4: Grid, Domain, and Runtime}

\subsection{Why resolution matters}
A finite-difference answer depends on \(\Delta S\), \(\Delta\tau\), and
\(S_{\max}\). Refining the grid reduces discretization error but increases
matrix size and the number of time steps. Boundary locations are especially
grid-limited because the exercise transition is first detected between
neighboring spot nodes.

The grid-240 solution is used as an internal numerical reference. It is not an
analytic truth surface. If \(Q_h\) is a quantity computed on grid \(h\), the
reported comparison has the form
\[
  \mathrm{RMSE}(Q_h,Q_{\mathrm{ref}})
  =
  \left(
    \frac{1}{|\Omega|}
    \sum_{(i,n)\in\Omega}
    \left(Q_{h,i}^n-Q_{\mathrm{ref},i}^n\right)^2
  \right)^{1/2},
\]
after interpolation where required. For Greeks, \(\Omega\) is restricted by
the strict mask.

\subsection{Domain truncation}
The true Black--Scholes spot domain is unbounded, whereas the computation stops
at \(S_{\max}\). If \(S_{\max}\) is too small, the imposed upper boundary can
affect economically relevant interior nodes. Comparing several multiples of
\(K\) tests this sensitivity but does not prove domain independence for every
future parameter regime.

\subsection{Convergence and runtime interpretation}
Price generally converges more smoothly than the free boundary, Delta, or
Gamma. Gamma converges most slowly because it is a second derivative and is
sensitive to nonsmooth initial data. Runtime rises with the number of time
steps, the number of spatial unknowns, and the PSOR sweeps required at each
step. The project reports runtime together with iteration counts so that speed
is not separated from solver convergence.

\begin{commonconfusion}
A small error relative to grid 240 means agreement with a finer internal
calculation. It does not mean agreement with an exact American-option formula,
because no such general closed form is being used as the reference.
\end{commonconfusion}

\section{Stages 5--6: Risk Surfaces and Dataset Governance}

\subsection{What a risk surface represents}
A risk surface evaluates price or a diagnostic over many spot, time, and model
parameter combinations rather than at one contract point. This reveals how
early exercise and sensitivities vary across regimes.

The application sweep separates:
\begin{itemize}
  \item American puts, whose exercise region is on the low-spot side;
  \item dividend-paying American calls, which may exercise on the high-spot
        side;
  \item no-dividend American call controls, which test the no-early-exercise
        theorem rather than provide dividend-call boundary evidence.
\end{itemize}
The verified application set contains 288 regimes: 144 puts, 108
dividend-paying calls, and 36 no-dividend controls. Thirty-two regimes carry
diagnostic failures and remain visible in the audit.

\subsection{Audit data versus clean labels}
Audit retention and model-label eligibility serve different purposes. A failed
or delicate regime remains scientifically important because it shows where the
solver diagnostics or domain assumptions need attention. It is excluded from
the clean machine-learning labels because training on a questionable target
would transfer numerical failure into the surrogate.

After controls and flagged regimes are excluded, the clean dataset contains
223 regimes and 3,264,943 samples. The 65 excluded regimes are recorded rather
than deleted from the project history.

\subsection{Why the split is by regime}
All spot-time nodes generated from one parameter tuple share the same option
surface. A pointwise random split could place neighboring points from the same
surface in both training and test sets. The model could then appear to
generalize while mostly interpolating within a regime it has already seen.

A regime-level split assigns the entire surface to one partition. The verified
split contains 156 training regimes, 33 validation regimes, and 34 test
regimes.

\begin{readerintuition}
Treat one parameter regime like one patient in a medical dataset: many
measurements from the same underlying case should not be divided between
training and testing.
\end{readerintuition}

\section{Stages 7--8: Direct Price and Positive Premium}

\subsection{Why low RMSE is insufficient}
A direct-price regressor learns
\[
  \widehat U_\theta=f_\theta(x).
\]
Standard squared-error training rewards average price accuracy but does not
encode \(\widehat U\geq\Phi\). A model can therefore achieve low RMSE while
predicting below the immediate payoff at some nodes.

A soft obstacle penalty may add
\[
  \lambda
  \left[\max\!\left(\Phi-\widehat U,0\right)\right]^2
\]
to the loss. This discourages violations but remains a trade-off with the data
loss. Finite \(\lambda\), finite optimization, and unseen inputs mean the
penalty is not a hard guarantee.

\subsection{Positive continuation-premium reconstruction}
Instead, predict a nonnegative premium:
\[
  \widehat C_\theta=\operatorname{softplus}(g_\theta(x))\geq0,
\]
\[
  \widehat U=\Phi+\widehat C.
\]
Therefore
\[
  \widehat U=\Phi+\widehat C,\quad \widehat C\geq0
  \Rightarrow
  \widehat U\geq\Phi.
\]
This removes one specific failure mode: payoff-obstacle violation.

\begin{commonconfusion}
Obstacle consistency is not full financial correctness. The construction does
not enforce monotonicity, convexity, boundary accuracy, Delta accuracy, or
Gamma accuracy. It also does not make the neural model a substitute for the
CN/PSOR benchmark.
\end{commonconfusion}

The selected positive-premium MLP has validation RMSE \(0.0200\), test RMSE
\(0.0218\), and zero obstacle violation on evaluated samples. The linear
premium sanity baseline has lower aggregate RMSE in this run, so the neural
model is not claimed to dominate all baselines.

\section{Stages 9--10: Boundary and Delta Diagnostics}

\subsection{Why boundary prediction is difficult}
The free boundary is an implicit transition derived from an entire option
surface. It changes orientation by option type, may be absent, and is limited
by the finite-difference grid. Boundary labels are also repeated over many
spot rows for the same regime-time pair, so a large point count does not imply
an equally large number of independent boundaries.

The boundary model is therefore a separate supervised diagnostic head:
\[
  \widehat B/K=h_\theta(\tau/T,r,q,\sigma,T,\text{option type}).
\]
It does not solve the obstacle PDE or reproduce PSOR iterations. The pooled
boundary MLP test RMSE is \(0.3513\) and its test MAE is \(0.2087\), making it
the weakest of the three neural components.

\subsection{Why price autograd Delta requires validation}
Let
\[
  y=\frac{U}{K},
  \qquad
  z=\log(S/K).
\]
Because \(U=Ky\) and
\[
  \frac{\partial z}{\partial S}=\frac{1}{S},
\]
the chain rule gives
\[
  \Delta
  =
  \frac{\partial U}{\partial S}
  =
  K\frac{\partial y}{\partial z}\frac{1}{S}
  =
  \frac{\partial (U/K)}{\partial \log(S/K)}
  \frac{1}{S/K}.
\]
A differentiable price network supplies this derivative computationally, but
price-loss accuracy does not directly train the derivative. Small level errors
can still create material slope errors.

The price-model autograd Delta test RMSE is \(0.0977\), which motivates a
separate supervised head. Its output map imposes
\[
  \Delta_{\mathrm{put}}\in[-1,0],
  \qquad
  \Delta_{\mathrm{call}}\in[0,1].
\]
The bounded Delta MLP has strict-mask test RMSE \(0.0332\), test MAE \(0.0140\),
and zero bound violation on evaluated samples.

\begin{commonconfusion}
Bounding Delta limits its range; it does not prove that the predicted slope is
correct. It also gives no second-derivative control and therefore does not
support a Gamma surrogate.
\end{commonconfusion}

\section{Stage 11: Integrated Scenario Screening}

\subsection{Why the workflow has three heads}
The integrated workflow keeps separate mathematical responsibilities:
\begin{enumerate}
  \item the positive-premium head estimates price while enforcing the payoff
        obstacle;
  \item the boundary head estimates a PSOR-derived transition as a diagnostic;
  \item the bounded Delta head estimates a first derivative within
        option-specific ranges.
\end{enumerate}
Combining their outputs in one interface does not turn them into one universal
model. Each head has separate labels, losses, metrics, and limitations.

\subsection{Scenario screening, not replacement}
Scenario screening means generating rapid preliminary estimates for cases
inside the tested regime envelope. Fresh CN/PSOR remains required for new
regimes, final label generation, and audit-quality decisions. The integrated
experiment evaluates six deterministic held-out regimes: three puts and three
dividend-paying calls.

Across those cases, the verified mean price RMSE is \(0.0092\), mean boundary
MAE is \(0.1539\), and mean strict-mask Delta RMSE is \(0.0300\). Obstacle and
Delta-bound violations are zero on the audited cases. Six cases, all from the
designed test envelope, are not enough to establish universal generalization.

\begin{readerintuition}
The workflow is a fast screening layer followed by a numerical audit, much as
a preliminary model can prioritize scenarios before a higher-fidelity
simulation is run.
\end{readerintuition}

\subsection{Overall Failure Modes and Limitations}
The expanded workflow retains the original conservative boundaries:
\begin{itemize}
  \item grid 240 is an internal reference, not analytic truth;
  \item Rannacher smoothing is not active;
  \item boundary estimation remains weaker than price and Delta estimation;
  \item strict masks exclude the most delicate Greek regions;
  \item the evidence does not support a Gamma surrogate;
  \item the neural workflow does not replace CN/PSOR;
  \item six integrated audit regimes do not establish broad generalization.
\end{itemize}

\subsection{Conclusion and Next Step}
The project demonstrates a benchmark-first learning workflow: formulate the
obstacle problem, validate CN/PSOR, attach diagnostics, govern labels, then
evaluate modular neural approximations under explicit claim limits. Stronger
claims would require broader independent regimes, higher-resolution numerical
references, and additional derivative validation.
